\documentclass[11pt]{article}
\usepackage[margin=1in]{geometry}
\usepackage{amsmath,amssymb,amsthm}
\usepackage[colorlinks=true,urlcolor=blue,linkcolor=blue,citecolor=blue]{hyperref}
\newtheorem{theorem}{Theorem}
\newcommand{\Med}{\operatorname{Med}}
\newcommand{\Prob}{\mathbb P}
\newcommand{\E}{\mathbb E}
\title{A proof of the generalized hyperbolic\\median lower bound}
\author{Principia Math in Progress}
\date{12 September 2026}
\begin{document}
\maketitle

\paragraph{Task 1: Resolve MathDB 345412.}
We prove Gaunt--Merkle's Conjecture 3.7 \cite{GM}, including the
variance-gamma boundary \(\delta=0\).
The proof uses the normal mixture representation, reciprocal pairing
of the mixing variable, and Cauchy--Schwarz.

\begin{theorem}\label{thm:main}
Let \(X\sim GH(\lambda,\alpha,\beta,\delta,0)\), where
\(\lambda>\tfrac12\), \(\beta>0\), \(\delta\geq0\), and
\(\gamma=\sqrt{\alpha^2-\beta^2}>0\). Then
\[
 \Med(X)>
 \frac{\beta}{\gamma^2}
 \left(\lambda-\frac12+
 \sqrt{\left(\lambda-\frac12\right)^2+\delta^2\gamma^2}\right).
\]
\end{theorem}

\paragraph{The mixing law and its reciprocal center.}
Let \(N\) be a standard normal variable independent of \(W\), and write
\begin{equation}\label{eq:mixture}
 X\ \stackrel d=\ \beta W+\sqrt W\,N.
\end{equation}
For \(\delta>0\), the generalized inverse Gaussian mixing density is
\begin{equation}\label{eq:gig}
 p(w)=C w^{\lambda-1}
       \exp\left\{-\frac12\left(\gamma^2w+\frac{\delta^2}{w}\right)\right\},
 \qquad w>0,\quad
 C=\frac{(\gamma/\delta)^\lambda}{2K_\lambda(\delta\gamma)}.
\end{equation}
For \(\delta=0\), use the gamma law with shape \(\lambda\) and rate
\(\gamma^2/2\). Integration of the conditional normal density in
\eqref{eq:mixture} gives precisely the GH density of \cite[eq.\ (1.1)]{GM}:
the exponent inside the mixing integral is
\[
 \beta x-\frac12\left((\gamma^2+\beta^2)w+
                            \frac{\delta^2+x^2}{w}\right),
\]
and \(\gamma^2+\beta^2=\alpha^2\).
Thus the representation uses exactly the source's parametrization.

Put \(k=\lambda-\tfrac12>0\), and define
\begin{equation}\label{eq:center}
 w_0=\frac{k+\sqrt{k^2+\delta^2\gamma^2}}{\gamma^2}>0,
 \qquad
 \gamma^2w_0-\frac{\delta^2}{w_0}=2k.
\end{equation}
We will show \(\Prob(X>\beta w_0)>1/2\).
Since the mixture has a continuous, strictly increasing distribution
function, this is equivalent to the theorem.

\paragraph{A half-moment comparison.}
We first establish
\begin{equation}\label{eq:moment}
 \E W^{1/2}-w_0\E W^{-1/2}\ \geq\ 0,
\end{equation}
with strict inequality for \(\delta>0\).
In that case all real moments of \(W\) are finite. Set
\(A=\E W^{1/2}\), \(B=\E W^{-1/2}\), and \(D=\E W^{-3/2}\).
Integration of the derivative of
\(w^{\lambda-1/2}\exp\{-(\gamma^2w+\delta^2/w)/2\}\)
over \((0,\infty)\), whose two boundary values vanish, gives
\begin{equation}\label{eq:recurrence}
 \gamma^2 A=2kB+\delta^2D.
\end{equation}
Strict Cauchy--Schwarz applied to \(W^{1/4}\) and \(W^{-3/4}\) gives
\(B^2<AD\); equality would require the nondegenerate variable \(W\)
to be constant. With \(R=A/B>0\), \eqref{eq:recurrence} implies
\[
 \gamma^2R^2-2kR
   =\delta^2\frac{AD}{B^2}>\delta^2.
\]
The quadratic has the unique positive root \(w_0\), so \(R>w_0\).
If \(\delta=0\), the gamma moment formula instead gives
\[
 \frac{\E W^{1/2}}{\E W^{-1/2}}
 =\frac{2}{\gamma^2}
     \frac{\Gamma(\lambda+1/2)}{\Gamma(\lambda-1/2)}
 =\frac{2k}{\gamma^2}=w_0.
\]
The negative half-moment is finite because \(\lambda>1/2\).
No negative three-half moment is required in this boundary case.

\paragraph{One change of sign after reciprocal pairing.}
Let \(U=\log(W/w_0)\), with density \(h\).
For both mixing laws above,
\[
 h(u)=C_0\exp\left\{\lambda u-
       \frac12\left(\gamma^2w_0e^u+
                       \frac{\delta^2}{w_0}e^{-u}\right)\right\},
 \qquad u\in\mathbb R,\quad C_0>0.
\]
By \eqref{eq:center},
\begin{equation}\label{eq:reflection}
 \log\frac{h(u)}{h(-u)}
      =2\lambda u-2k\sinh u.
\end{equation}
The function \(\sinh u/u\) increases strictly from \(1\) to infinity
on \(u>0\): its derivative has numerator \(u\cosh u-\sinh u>0\).
Since \(\lambda/k>1\), there is a unique \(u_*>0\) such that
\(\sinh u_*/u_*=\lambda/k\).
Consequently, with \(\Delta(u)=h(u)-h(-u)\),
\begin{equation}\label{eq:sign}
 \Delta(u)>0\quad(0<u<u_*),\qquad
 \Delta(u)<0\quad(u>u_*).
\end{equation}
Writing \(s(u)=\sinh(u/2)\), the half-moment comparison becomes
\begin{equation}\label{eq:weighted}
 \int_0^\infty s(u)\Delta(u)\,du
 =\E\sinh(U/2)
 =\frac{\E W^{1/2}-w_0\E W^{-1/2}}{2\sqrt{w_0}}
 \geq0.
\end{equation}
All terms are integrable by the moment conditions already checked.

\paragraph{The normal tail preserves the needed sign.}
Let \(\Phi\) and \(\varphi\) be the standard normal distribution function
and density, and let \(c=2\beta\sqrt{w_0}>0\).
Conditional on \(W=w_0e^u\),
\[
 \Prob(X>\beta w_0\mid U=u)
 =\Phi\left(\beta\,\frac{W-w_0}{\sqrt W}\right)
 =\Phi(c\sinh(u/2)).
\]
For \(s>0\), define
\[
 \rho(s)=\frac{\Phi(cs)-1/2}{s}
        =c\int_0^1\varphi(csv)\,dv.
\]
This function is positive and strictly decreasing. Pairing \(u\) with
\(-u\) therefore gives
\begin{align}
 \Prob(X>\beta w_0)-\frac12
 &=\int_0^\infty \rho(s(u))s(u)\Delta(u)\,du \notag\\
 &=\rho(s(u_*))\int_0^\infty s(u)\Delta(u)\,du \notag\\
 &\quad+\int_0^\infty
       [\rho(s(u))-\rho(s(u_*))]s(u)\Delta(u)\,du.            \label{eq:conclusion}
\end{align}
The first term is nonnegative by \eqref{eq:weighted}.
The integrand in the second term is strictly positive on both
\((0,u_*)\) and \((u_*,\infty)\), by \eqref{eq:sign} and the strict
decrease of \(\rho\). Both integrals are finite, since \(\rho\) is bounded
and \(\int s|\Delta|<\infty\).
Thus the right-hand side is strictly positive. This proves
Theorem \ref{thm:main}, including strictness when \(\delta=0\), where
the first term vanishes.

\paragraph{An extension of the parameter range.}
For \(\delta>0\), the same bound actually holds for every real \(\lambda\).
If \(\lambda\leq1/2\), keep \(k=\lambda-1/2\) and the positive root
\eqref{eq:center}. Then, for \(u>0\), \eqref{eq:reflection} equals
\[
 u+(1-2\lambda)(\sinh u-u)>0.
\]
In this case \(\Delta(u)>0\) everywhere, and the first paired integral
in \eqref{eq:conclusion} is positive directly. The half-moment argument
and a sign-change point are unnecessary for this extension.

\paragraph{Verification and status.}
The companion verifier checks the exact reflection and moment algebra,
computes the normal-mixture distribution function at the proposed bound,
and cross-checks special cases using the direct GH density and the
asymmetric Laplace law. These calculations supplement the analytic proof.
The live MathDB entry remained open with empty claim fields.
The current source, its indexed published statement, and subsequent
public literature were checked; no matching full claim was located.
The dated audit records the scope of those searches.

\paragraph{Task 2: Find the best bound uniform in the skew parameter.}
For fixed mixing parameters, determine the infimum of
\(\Med(\beta W+\sqrt W N)/\beta\) over \(\beta>0\), and compare it
with the reciprocal center \(w_0\) and the half-moment ratio.

\paragraph{Task 3: Quantify the gap above the lower bound.}
Derive explicit positive lower estimates for
\(\Med(X)-\beta w_0\), including uniform asymptotic formulas as
\(\delta\gamma\) tends to zero or infinity.

\paragraph{Task 4: Extend reciprocal pairing to other mixing laws.}
Characterize positive mixing densities for which the reflected
log-density has the sign pattern \eqref{eq:sign}, and develop
median bounds from appropriate half-moment inequalities.

\paragraph{Task 5: Replace the normal mixing noise.}
Determine which symmetric noise densities allow the same argument,
including the strict decrease of the corresponding tail-average
function \(\rho\), and identify the resulting mixture families.

\begin{thebibliography}{9}
\bibitem{GM}
R.\ E.\ Gaunt and M.\ Merkle,
On bounds for the mode and median of the generalized hyperbolic and
related distributions,
\emph{J.\ Math.\ Anal.\ Appl.} \textbf{493} (2021), 124508.
\href{https://doi.org/10.1016/j.jmaa.2020.124508}{doi:10.1016/j.jmaa.2020.124508}.
\href{https://arxiv.org/abs/2002.01884}{arXiv:2002.01884v4},
Conjecture 3.7.
\end{thebibliography}
\end{document}
